\documentclass[12pt]{article}
\usepackage[margin=1in]{geometry}
\usepackage{setspace}
\usepackage{parskip}
\usepackage{enumitem}
\usepackage{titlesec}
\usepackage{xcolor}
\usepackage{mdframed}
\usepackage{hyperref}
\usepackage{amsmath}

% -------------------------------------------------------
% Styling
% -------------------------------------------------------
% Referee 1 � blue
\definecolor{r1color}{RGB}{30, 80, 150}
\definecolor{r1bg}{RGB}{240, 245, 255}
% Referee 2 � teal
\definecolor{r2color}{RGB}{0, 110, 100}
\definecolor{r2bg}{RGB}{235, 248, 245}

\newmdenv[
  backgroundcolor=r1bg,
  linecolor=r1color,
  linewidth=1.5pt,
  leftline=true, rightline=false, topline=false, bottomline=false,
  innerleftmargin=10pt, innerrightmargin=10pt,
  innertopmargin=6pt, innerbottommargin=6pt,
  skipabove=8pt, skipbelow=4pt
]{r1quote}

\newmdenv[
  backgroundcolor=r2bg,
  linecolor=r2color,
  linewidth=1.5pt,
  leftline=true, rightline=false, topline=false, bottomline=false,
  innerleftmargin=10pt, innerrightmargin=10pt,
  innertopmargin=6pt, innerbottommargin=6pt,
  skipabove=8pt, skipbelow=4pt
]{r2quote}

\titleformat{\section}{\large\bfseries\color{currcolor}}{}{0em}{}[\titlerule]
\titleformat{\subsection}{\normalsize\bfseries}{}{0em}{}

% -------------------------------------------------------
% Document metadata
% -------------------------------------------------------
\newcommand{\manuscriptid}{YGAME-D-25-00181}
\newcommand{\papertitle}{Preferences for Warning Signal Quality: Experimental Evidence}
\newcommand{\journalname}{Games and Economic Behavior}
\newcommand{\submissiondate}{\today}

% -------------------------------------------------------
\begin{document}

\begin{center}
  {\LARGE\bfseries Response to Referee Reports}\\[6pt]
  {\large Manuscript \manuscriptid}\\[4pt]
  {\normalsize \textit{\papertitle}}\\[4pt]
  {\normalsize \journalname\ $\cdot$ \submissiondate}
\end{center}

\vspace{1em}
\hrule
\vspace{1em}

We thank the associate editor and both referees for their thorough and constructive
reviews. The revised manuscript is stronger for it. Below we address each point in
turn. Referee comments are shown in shaded boxes ({\color{r1color}blue} for
Referee~1, {\color{r2color}teal} for Referee~2); our responses follow in plain text.
%
%% ===============================================================
%\colorlet{currcolor}{black}
%% ===============================================================
%
%\tableofcontents
%\newpage
%
%% ===============================================================
%\colorlet{currcolor}{black}
%\section{Associate Editor}
%% ===============================================================
%
%The Associate Editor has no additional comments.
%
%% ===============================================================
%\newpage
\colorlet{currcolor}{r1color}
\section{Referee 1}
% ===============================================================

We thank Referee~1 for the careful reading of both the original submission and the
revision.

%----------------------------------------------------------------
\subsection{R1.1\quad Major Comment: Incorrect and Inconsistent Definitions of FP and FN Costs}

\begin{r1quote}
The paper now emphasizes the distinction between false-positive/false-negative rates and
false-positive/false-negative costs. This is a useful clarification. However, the definitions
given in the model and in Section 5.2 are not correct. The paper defines FP and FN costs on
page 6 as $\pi(1 - P_{01})c$ and $\pi P_{01}L$, respectively, and states that WTP is equally affected by
both. There are two problems.

First, the FP cost definition is incorrect. The expression $\pi(1-P_{01})c$ is the expected cost
of true positive protection (correctly protecting in the adverse state), not the cost of false
positives. The FP cost should be the expected cost of unnecessary protection when there is
no adverse event: $(1 - \pi)P_{10}c$. This is consistent with equation (4), where the derivative of
risk-neutral WTP with respect to the FP rate is $-(1 - \pi)c$.

Second, there is a subtlety in the FN cost definition. The paper defines it as $\pi P_{01}L$, which
is the gross expected loss from missed detections. An alternative definition is $\pi P_{01}(L - c)$,
which captures the incremental cost relative to the correct action: under a true positive
signal the agent would have protected and paid $c$, so the net cost of missing the adverse
state is $L - c$, not $L$. This alternative is more consistent with equation (5), where the
derivative of WTP with respect to the FN rate is $-\pi(L - c)$.

This matters for Table 5 and the main regression. The dependent variable is WTP minus
the risk-neutral value, and the independent variables are labeled FP and FN costs. The
paper should explicitly define the variables used in the regressions, show how they map
into the theoretical WTP formula (equation 3), and verify that the definitions do not affect
the reported results.
\end{r1quote}

\par\smallskip\noindent\textbf{Response:} [Your response here.]

%----------------------------------------------------------------
\subsection{R1.2\quad Major Comment: The Interpretation of Table 5 and Figure 5 Should Be Stated More Precisely}

\begin{r1quote}
I found that the interpretation and conclusion of the paper are not clear and somewhat
misleading. For instance, the abstract states that subjects ``disproportionately reduce their
demand for signals with high false negative rates for rare events.'' This characterization is
imprecise and risks confusion with what the data actually show in subsection 5.2 (conditional on the results are correct).

Under the null hypothesis that subjects are risk-neutral expected-utility maximizers,
the benchmark WTP already incorporates FP and FN costs, so all coefficients in the discrepancy regression (Table 5) should be zero. Positive coefficients indicate deviations from the
benchmark: as error costs rise, subjects' WTP does not fall enough. The actual finding of
the regression table is that subjects selectively under-react to the cost type that the theory
identifies as dominant: for low priors, the coefficient on FP costs is large and significant
while the coefficient on FN costs is near zero; for high priors, the pattern reverses, with FN
costs generating the significant positive coefficient. The theory (equations 4--5) predicts
that FP costs should be the primary concern for rare events, since $(1 - \pi)$ is large and most
positive signals will be false alarms, while FN costs should dominate for frequent events,
since $\pi$ is large and missed detections become costly. This implies that subjects under-react
to precisely the cost type that theory says should matter most:
\begin{itemize}
  \item Low priors $\Rightarrow$ FP costs should dominate $\Rightarrow$ subjects fail to adjust for FP costs
  \item High priors $\Rightarrow$ FN costs should dominate $\Rightarrow$ subjects fail to adjust for FN costs
\end{itemize}

Figure 5 corroborates this from a different angle: when raw WTP is regressed on FP and
FN rates by prior, the empirical sensitivity coefficients are roughly equal and flat across
priors, whereas theory predicts sharp divergence. The finding is not that subjects overpenalize FN rates for rare events in absolute terms, but rather that they fail to differentially
weight FP and FN rates as the prior shifts, producing relative over-sensitivity to FN (undersensitivity to FP) for low priors and the reverse for high priors. I think this is stronger than
saying simply that subjects are ``asymmetrically under-responsive by prior,'' but weaker
than saying that subjects generally fail to distinguish FP and FN errors. The evidence shows
that subjects do not process FP and FN rates in the benchmark Bayesian/EU way. It does
not uniquely identify the cognitive mechanism.

The abstract and conclusion should be rewritten to state this more precisely. The paper
should also be consistent about whether claims refer to error rates, expected costs, or
deviations from the risk-neutral benchmark. The paper would also benefit from using ``false
alarm'' and ``missed detection'' alongside the FP/FN notation, especially when discussing
costs. These terms make it immediately clear why low-prior environments make false alarms
more important and high-prior environments make missed detections more important.
\end{r1quote}

\par\smallskip\noindent\textbf{Response:} XXX Rewrite such that we do not claim subjects don't know FP/FN errors. "`Subjects to account the effect of priors on sensitivity of expected costs to FP/FN rates"'

%----------------------------------------------------------------
\subsection{R1.3\quad Major Comment: The Mechanism Claim Should Be Softened}

\begin{r1quote}
The discussion in Section 5.3 is much improved. The authors now consider risk preferences,
probability weighting, anchoring, non-instrumental information, and anticipation of biased
protection decisions. The additional BE and IP analyses are also useful. However, I remain
unconvinced that the data uniquely identify a failure to distinguish false-positive and falsenegative errors. The evidence is interesting: some subjects report using the proportion of
dishonest gremlins as a decision rule, and the BE/IP analyses show anomalous responses to
FP and FN rates even when those rates should not affect the posterior given the observed
signal (Tables 6--7).

However, the mechanism is not cleanly identified. The same pattern could reflect a
broader total-error-rate heuristic, a ``dishonest gremlin'' heuristic induced by the experimental framing, probability weighting, noisy Bayesian reasoning, task complexity, or some
combination. The authors themselves acknowledge that probability weighting may play a
role for low priors. More importantly, the leap from ``subjects react to irrelevant error rates
in BE and IP tasks'' to ``subjects conflate FP and FN when valuing signals'' requires more
thorough discussion and the implicit assumption that the same cognitive failure operates
across tasks.

A more defensible conclusion would be: under this experimental representation, many
subjects appear to use a simplified total-error or total-dishonesty heuristic rather than
separately processing how FP and FN rates interact with the prior probability. This interpretation fits the subject explanations and the BE/IP evidence better than the stronger
claim that subjects generally cannot distinguish FP and FN errors. It is still an interesting
finding.

Relatedly, Section 4 should not be framed only as reassurance that subjects understood
the task. The BP, IP, and BE patterns are part of the behavioral phenomenon. The data
suggest that subjects respond directionally to incentives and information, but also display
substantial heterogeneity, mistakes, and simplified reasoning. That framing would make
the transition to the WTP results more natural.
\end{r1quote}

\par\smallskip\noindent\textbf{Response:} XXX Rewrite to weaken the claims, follow suggestion XXX

%----------------------------------------------------------------
\subsection{R1.4\quad Minor Comment: Proposition 2}

\begin{r1quote}
Ensure the wording explicitly states that the proposition concerns the
ratio of sensitivities $(\partial b/\partial P_{01})/(\partial b/\partial P_{10})$, not absolute levels. A brief economic interpretation would help: compared to risk-neutral agents, risk-averse agents are relatively more
intolerant of missed detections (false negatives) but more forgiving of false alarms (false
positives).
\end{r1quote}

\par\smallskip\noindent\textbf{Response:} [Your response here.]

%----------------------------------------------------------------
\subsection{R1.5\quad Minor Comment: Appendix B Proof Numbering}

\begin{r1quote}
The main text contains two propositions, but the appendix has three proofs (B.1--B.3). The proof for Proposition 2 (ratio of sensitivities) appears
in B.3. Section B.2 appears to prove a result for imprudent agents that no longer corresponds
to a proposition in the main text. B.2 should be dropped or the main text should be updated
to match.
\end{r1quote}

\par\smallskip\noindent\textbf{Response:} We have now dropped the old Section B.2 that had a proof that is no longer relevant. Thank you for pointing this out.

%----------------------------------------------------------------
\subsection{R1.6\quad Minor Comment: Appendix B.3 Algebra Error}

\begin{r1quote}
In equations (13) and (14), the denominator of the last
expression is $(1 - \pi)L$. It should be $(1 - \pi)c$. This can be verified by simplifying the middle
expression:
\[
\frac{\pi}{1 - \pi} \cdot \frac{u'(\xi_1)(L - c)}{u'(\xi_2) c} = \frac{\pi(L - c)}{(1 - \pi) c} \cdot \frac{u'(\xi_1)}{u'(\xi_2)}.
\]
\end{r1quote}

\par\smallskip\noindent\textbf{Response:} We thank the referee for catching this typo. We have corrected the
denominator in equations~(13) and~(14) from $(1-\pi)L$ to $(1-\pi)c$. The stated result
is unaffected.

%----------------------------------------------------------------
\subsection{R1.7\quad Minor Comment: Cross-Reference Error}

\begin{r1quote}
Page 26 references ``Appendix Table A10,'' but the appendix table
with matching content (WTP deviations with IP bias controls) is numbered A8. The table
numbering appears to be off.
\end{r1quote}

\par\smallskip\noindent\textbf{Response:} [Your response here.]

%----------------------------------------------------------------
\subsection{R1.8\quad Minor Comment: Typos}

\begin{r1quote}
\begin{itemize}[noitemsep]
  \item Page 7: ``Equations 6-s 7'' $\to$ ``Equations 6--7.''
  \item Page 15, footnote 10: missing comma between the two parameter estimates ($\hat{\pi} = 0.43$ $\hat{\theta} = 0.25$).
  \item Page 18: ``the value for for a prior of 0.2'' --- duplicated word.
  \item Page 27: ``the observed sensitivity patterns emerges'' $\to$ ``emerge.''
  \item Page 27: ``characteristcs'' $\to$ ``characteristics.''
\end{itemize}
\end{r1quote}

\par\smallskip\noindent\textbf{Response:} All five typos have been corrected. We thank the referee for the
careful reading.

% ===============================================================
\newpage
\colorlet{currcolor}{r2color}
\section{Referee 2}
% ===============================================================

We thank Referee~2 for the detailed and constructive comments on the revision.

%----------------------------------------------------------------
\subsection{R2.1\quad Treatments and Treatment Assignment}

\begin{r2quote}
\textbf{(a)} While the new ``Treatments and Treatment Assignment'' subsection on page 10 provides additional detail on how treatments are assigned, it remains incomplete. It
states that each subject is randomly allocated to a pair of priors but does not explain
how information structures are assigned.

\textbf{(b)} The text also notes that ``the first treatment was always the one with a fully truthful
information structure,'' which suggests that information structures are not randomly
assigned across subjects. For transparency and to avoid confusion, the authors
should provide at least a general description and justification of how information
structures are allocated.

\textbf{(c)} Because information structures do not appear to be randomly assigned, the resulting
sample sizes across treatment cells (i.e., prior $\times$ information structure combinations)
are likely to differ. I encourage the authors to report sample sizes by treatment cell.
Given the large number of cells ($4 \times 12 = 48$), a table in the appendix would likely
suffice.
\end{r2quote}

\par\smallskip\noindent\textbf{Response:} Thank you for the suggestion. We have now added the full lists of treatments and sequences in the appendix. We have also revised the description of the treatment assignment in Section 3 as follows:

\begin{quotation}

\noindent\textbf{Treatments and Treatment Assignment.} Table~XX summarizes our treatments for tasks other than BP. The treatments combine four priors with 12 different information structures for a total of 48 treatments. These treatments allow for a distribution of posteriors (for black and white signals) that spans a wide range of interior values and allows for a wide range of theoretical WTP values (Appendix Figure~XX). 

Based on these 48 treatments, we constructed 32 unique treatment sequences that were randomly assigned to subjects. Each sequence comprises six rounds of treatments that combined a pair of priors of either (0.1, 0.3) or (0.2, 0.5) with three information structures. Within a sequence, treatments were presented by prior; in half of the sequences, the first prior is larger than the second. For each prior, the first treatment was always the one with a fully truthful information structure. The order of priors and signals stays constant across tasks, but vary between sequences. Appendix Tables~A1 and~A2 present the full lists of treatments and sequences.

\end{quotation}

%----------------------------------------------------------------
\subsection{R2.2\quad Description and Groupings in Tables 2, 3, and A4}

\begin{r2quote}
\textbf{(a)} On page 12, the text states that ``Table 2 presents the average protection decisions
by prior and signal type,'' but Table 2 does not, in fact, report descriptive statistics by prior. Instead, it groups observations by signal value (white/black) and by
information-structure type (truthful, FP-only, FN-only, both errors). The same mismatch appears in the caption of Appendix Table A4, which refers to ``signal type
and prior,'' although the table is actually organized by signal value and informationstructure type.

\textbf{(b)} Relatedly, the authors could briefly justify why Table 2 groups observations by signal value and information structure while pooling across priors. A single sentence
would suffice. Providing this explanation would improve transparency and clarity
for readers, especially since the same grouping is used in Tables 3 and A4.
\end{r2quote}

\par\smallskip\noindent\textbf{Response:} Typo, we moved the tables to the appendix. Adjust the text in the paper. 

%----------------------------------------------------------------
\subsection{R2.3\quad Claim on Page 16 Regarding the Introduction of the FN Rate (Row 1 vs.\ 2)}

\begin{r2quote}
\textbf{(a)} The paragraph below Table 3 states that introducing the FP rate (row 1 vs.\ 3 of
Table 3) magnifies the upward updating bias for white signals---which is correct: the
Bayesian posterior remains at 0, while the mean updating error increases from 0.048
to 0.210. The text then claims that ``we find a similar effect for the introduction of
the FN rate (row 1 vs.\ 2).'' However, comparing row 1 and row 2, the mean updating
error does not appear to change significantly (0.048 vs.\ 0.051).

\textbf{(b)} Moreover, introducing the FN rate changes the Bayesian posterior for the white
signal (from 0.000 to 0.119), so the comparison between row 1 and row 2 does not
hold the posterior fixed---unlike the comparison between row 1 and row 3.
\end{r2quote}

\par\smallskip\noindent\textbf{Response:} Correct the text

%----------------------------------------------------------------
\subsection{R2.4\quad Section 4's Closing Summary Understates the Evidence in Tables 2 and 3}

\begin{r2quote}
\textbf{(a)} The paragraph just before Section 5 (page 17) concludes that ``All of this suggests
that subjects understand the experimental environment well enough to allow us to
draw reasonable inferences from our WTP elicitation.'' While I agree that subjects
appear to understand the environment sufficiently well, this framing understates
several important findings from Section 4.

\textbf{(b)} For example, Tables 2 and 3 show that increasing the FP (FN) rate affects protection
decisions and updating errors even in cases where the rate should not influence the
Bayesian posterior (e.g., FP for a white signal in Table 3, row 1 vs.\ 3). These
patterns provide suggestive evidence of the FP/FN-confusion mechanism. I would
encourage the authors to add a sentence foreshadowing that the descriptive analysis
already reveals patterns that are later revisited and explained in Section 5.3. This
would make the main mechanism visible from Section 4 onward, rather than only at
the end of the paper.
\end{r2quote}

\par\smallskip\noindent\textbf{Response:} Add it back in foreshadowing this result.

%----------------------------------------------------------------
\subsection{R2.5\quad Hypothesis 1 and Result 1 Overstate What the t-Tests in Table 4 Establish}

\begin{r2quote}
\textbf{(a)} Hypothesis 1 (page 17) states that ``subjects' WTPs for signals are equal to their
value for risk-neutral agents,'' and Result 1 concludes that ``on average, there are
no significant discrepancies between the empirical WTP and the risk-neutral benchmark.'' However, the supporting analysis relies on a t-test of the mean of $\Delta b_{is} =
b_{is} - b^*_s$. Failing to reject $E[\Delta b_{is}] = 0$ is a much weaker result than establishing
$\Delta b_{is} = 0$, as large but offsetting over- and under-valuations across subjects and
treatments may cancel out in the mean. Given the likely presence of individual
heterogeneity, this distinction is important. I suggest one of the following revisions:
\begin{enumerate}[label=\roman*.]
  \item Restate Hypothesis 1 and Result 1 more precisely, for example: ``subjects' WTPs
are not systematically above or below the risk-neutral benchmark on average,''
which is itself an informative finding; or
  \item Test the stronger claim directly, for instance by examining whether the absolute
deviation $|\Delta b_{is}|$ is significantly different from zero.
\end{enumerate}
\end{r2quote}

\par\smallskip\noindent\textbf{Response:} Do (i), explain why (ii) is not necessarily doable

%----------------------------------------------------------------
\subsection{R2.6\quad (Optional) Reorganize Section 5.3 to Give Greater Prominence to the FP/FN-Confusion Mechanism}

\begin{r2quote}
\textbf{(a)} Given that the FP/FN-confusion mechanism is the primary mechanism, it might be
preferable to present it as the first candidate explanation rather than the last, and
then proceed to evaluate and rule out alternative mechanisms.
\end{r2quote}

\par\smallskip\noindent\textbf{Response:} Reorganize to follow suggestion 

%----------------------------------------------------------------
\subsection{R2.7\quad Risk-Preference Discussion (pp.\ 23--24)}

\begin{r2quote}
In the risk-preference discussion (pp.\ 23--24), the authors could also note that the main
effects of FP and FN costs remain significant, rather than focusing solely on the insignificance of the interaction terms.
\end{r2quote}

\par\smallskip\noindent\textbf{Response:} Follow

%----------------------------------------------------------------
\subsection{R2.8\quad Probability Weighting Discussion (pp.\ 24--25)}

\begin{r2quote}
\textbf{(a)} In equation 9 and the definition of $\Lambda$, the numerator should be $P_{01}$ instead of $P_{00}$?

\textbf{(b)} I appreciate the authors' careful argument with rank-dependent utility! In my previous report, I was meant to refer to the probability weighting function from prospect
theory (Kahneman and Tversky, 1979). But the arguments turn out to be quite
similar because the particular payoff set used in the experiment.

\textbf{(c)} The flow of the argument is clear, and the observation that ``$\pi + (1 - \pi)P_{10} \leq
0.75$, meaning that the value is still likely to fall in the mid-range of an S-shaped
weighting function,'' is insightful. However, this depends on the rotation point of the
probability weighting function. Theoretically, if the rotation point lies far to the left,
then $\pi + (1 - \pi)P_{10} \leq 0.75$ might still fall within the steeper region of the weighting
function (such that the denominator is large and $\Lambda$ is small, implying lower relative
sensitivity). It would be helpful if the authors could provide evidence for ``middle''
rotation points in existing studies (e.g., Gonzalez and Wu (1999), as cited). That
said, since the shape of the weighting function is not restricted, this possibility may
be too flexible to rule out. In any case, it would be helpful to discuss this issue more
thoroughly.
\end{r2quote}

\par\smallskip\noindent\textbf{Response:} [Your response here.]

%----------------------------------------------------------------
\subsection{R2.9\quad The ``Anticipation of Bias in Protection Decision''}

\begin{r2quote}
\textbf{(a)} Table A10 should be Table A8; the caption of Table A8 should refer to WTP deviations rather than simply WTP.

\textbf{(b)} More importantly, I am not entirely convinced that the two newly introduced variables adequately test the mechanism of ``anticipation of bias in protection decisions.''
If I understand correctly, $\Delta\text{prot}$ measures the extra cost incurred by the subject relative to the optimal benchmark, while $\Delta\text{loss}$ measures the additional loss incurred
relative to the same benchmark. The coefficients in Table A8 are negative, suggesting that greater mistakes (in terms of cost or loss) reduce WTP deviations. In other
words, subjects who perform worse in the IP task appear to be more accurate in
their WTP, which may indicate that they account for their own mistakes.

\textbf{(c)} Given that the coefficients on the two new variables are statistically significant, the
conclusion that ``subjects' biases in how they will utilize signals when facing an IP
task do not explain the observed pattern of sensitivities to FP and FN rates'' seems
premature. A more appropriate interpretation is that the coefficients on FP and FN
do not appear to change after introducing these two variables.

\textbf{(d)} However, a more informative regression specification would include interaction terms
between the mistake variables and FP/FN rates. The intuition is that, if anticipation
of bias is the primary channel, then mistakes in protection decisions should mediate
the relationship between WTP deviations and FP/FN. That is, if subjects expect
to make more mistakes in the IP task at higher FP/FN rates, they may reduce
their WTP, leading to larger deviations. Along similar lines, it may be preferable
to distinguish mistakes in BP and IP directly, rather than separating them into cost
and loss components.

\textbf{(e)} Finally, even if anticipation of bias is an important (or primary) channel, it does not
undermine the FP/FN-confusion mechanism; if anything, it reinforces it. Tables 2
and 7 already show that subjects' IP protection decisions respond to FP and FN rates
in ways consistent with FP/FN confusion. For instance, subjects choose to protect
upon receiving a white signal even when there are no false negatives but only false
positives (Table 7 provides related evidence). Anticipation of bias suggests that
this confusion arises not only at the WTP stage but also during the IP stage, with
such anticipation leading to larger deviations of WTP from the risk-neutral optimal
benchmark.
\end{r2quote}

\par\smallskip\noindent\textbf{Response:} [Your response here.]

%----------------------------------------------------------------
\subsection{R2.10\quad Cross-Task Consistency Analysis}

\begin{r2quote}
The authors should consider incorporating the cross-task consistency analysis from the
response letter (Table 2 on page 13) into the paper, or at least including it in the appendix.
\end{r2quote}

\par\smallskip\noindent\textbf{Response:} Move it into the appendix

%----------------------------------------------------------------
\subsection{R2.11\quad (Optional) The ``Dishonest Gremlin'' Framing}

\begin{r2quote}
I understand the authors' argument that re-running the experiment with a different framing of the information structure might make pooling unjustifiable. However, the authors could consider adding a brief discussion of this issue (perhaps at the end of the paper). Even without the ``dishonest'' framing, individuals may still conflate the two types of errors (FP and FN). It would therefore be helpful for the authors to acknowledge and discuss this potential limitation.
\end{r2quote}

\par\smallskip\noindent\textbf{Response:} Add a discussion + a footnote why we didn't incorporate the neutral frame because of pooling concerns.

%----------------------------------------------------------------
\subsection{R2.12\quad Exposition}

\begin{r2quote}
\textbf{(a)} On page 11, the sentence ``this section also provides \ldots\ the theoretical predictions''
is a tad imprecise: the predictions refer to those of the risk-neutral benchmark,
not theoretical predictions in general. I would suggest ``\ldots\ from the theoretical
predictions of the risk-neutral benchmark'' (or similar) for clarity.

\textbf{(b)} Relabel ``Posterior'' as ``Bayesian posterior'' in Tables 2, 3, A2, and A4. It would
help the readers in case they mistake it as elicited posterior.

\textbf{(c)} In both Table 5 and Table A8, the boundary between the main regression panel and
the lower panel reporting the risk-averse and risk-loving subgroup tests is not easy
to spot at a glance. A horizontal rule separating the two would make the structure
of these tables much easier to spot.

\textbf{(d)} (Optional.) The appendix tables would be easier to navigate if the captions (or
notes) pointed the reader to the coefficient or column that matters for the corresponding discussion in the main text. A short closing sentence in each caption ---
e.g., ``comparing the coefficient of XYZ and ABC in column 2 suggests that \ldots'' ---
would substantially reduce the cost of cross-referencing.

\textbf{(e)} On page 21, the text refers to ``the systematic WTP biases reported above'' without
specifying which biases are meant. I would suggest the authors to be more precise
here.
\end{r2quote}

\par\smallskip\noindent\textbf{Response:} XXX FOLLOW AFTER MAIN ANALYSIS DONE XXX

%----------------------------------------------------------------
\subsection{R2.13\quad Typos}

\begin{r2quote}
\begin{itemize}[noitemsep]
  \item There is an extra ``$=$'' sign in equation (2) on page 5.
  \item On page 7, ``Equations 6-s 7'' should read ``Equations 6 and 7'' (or similar).
  \item Footnotes 10 and 11 are missing a full stop at the end.
  \item Footnote 13 (page 21): ``average these errors across all the decisions with identical
priors, false positive, and false negative rates''. Should ``identical'' be ``different''?
  \item In the last paragraph on page 26, Table A10 should be Table A8.
\end{itemize}
\end{r2quote}

\par\smallskip\noindent\textbf{Response:} All five typos have been corrected. We thank the referee for the
careful reading. XXX MAKE SURE WE DO THIS XXX

%----------------------------------------------------------------
\end{document}